\subsubsection{Oversampling}
\label{sec:semantics:grsync:theorems:oversampling}

Components are compiled into synchronous state machines, but their execution within a service
involves an additional optimization: computations are performed only when the inputs have changed.
The typing of component applications in services ensures that only \emph{reactive} components are
used (see \Cref{sec:semantics:grsync:react_typing}).
%
A reactive component emits events and uses memory and time-dependent expressions only when an event
has occurred or a signal has changed. This restriction is enforced by the reactive type system
described in \Cref{sec:semantics:grsync:react_typing}.
\Cref{thm:semantics:grsync:theorems:oversampling} states
that for components satisfying this type system, it is sufficient to propagate only input changes.
Therefore, additional evaluations of the components, referred to as \emph{oversampling} their
executions, do not alter the values of their outputs. This property enables a shift from periodic
to event-driven computations without risking inconsistency.

Similar to the previous theorems on components, the proof of
\Cref{thm:semantics:grsync:theorems:oversampling} proceeds by structural induction on the syntax of
\GRSync. Due to the presence of component calls within expressions, the proof is mutually recursive
with the proofs of the intermediate \Cref{%
    lem:semantics:grsync:theorems:oversampling:e,%
    lem:semantics:grsync:theorems:oversampling:pat,%
    lem:semantics:grsync:theorems:oversampling:epat,%
    lem:semantics:grsync:theorems:oversampling:eq%
}. These lemmas establish the oversampling property for each syntactic construct of \GRSync
(expressions, patterns, event patterns, and equations). This mutual recursion is well-founded
because, while expressions can contain component calls, the components themselves are not mutually
recursive, ensuring that the induction remains valid.

\begin{remark}
    Let $sty$ be a sampling type, and let $v$ and $v'$ be two values of type $sty$ (\eg if $sty$ is
    \intTY or \mayB{\intTY}, then the values are integers).
    %
    We use the notation $v' = \intoVal{sty}{v}{\bot}$ to express that the value $v'$ is an
    \emph{oversampling} of $v$.

    To understand this encoding, consider the two possible cases for $sty$:
    %
    \begin{itemize}
        \item If $sty$ is a signal type, then the definition of \smf{into\_val_{sty}}
              (\Cref{rule:semantics:grfrp:opsem:aux:intoVal:signal}) yields $v' = v$. This
              represents the absence of change for the signal, which is the expected behavior in the
              case of oversampling.

        \item If $sty$ is an event type, then the definition of \smf{into\_val_{sty}}
              (\Cref{rule:semantics:grfrp:opsem:aux:intoVal:event}) yields $v' = \none$. This
              represents the absence of the event, which is the expected behavior in the case of
              oversampling.
    \end{itemize}
\end{remark}

\begin{lemma}[Oversampling of \GRSync reactive expressions]
    \label{lem:semantics:grsync:theorems:oversampling:e}
    Let $e$ be an expression in a program \G and typing context \tyCtx, let $sty^+$ be types, let
    $\rType^+$ be reactive types, and let $\Theta$ be a reactive typing context such that:
    \begin{equation*}
        \exprTyped{e}{sty^+} \qquad \exprTypedReact{\Theta}{e}{\rType^+}.
    \end{equation*}
    For all states $s_1$, $s_2$, and $s_3$ corresponding to $e$, all contexts $\rho_1$ and $\rho_2$,
    and all values $v_1^+$ and $v_2^+$ such that:
    \begin{equation*}
        \exprOpsem{\rho_1}{s_1}{e}{s_2}{v_1^+} \qquad \exprOpsem{\rho_2}{s_2}{e}{s_3}{v_2^+},
    \end{equation*}
    we have:
    \begin{multline*}
        \left(\forall x \in \deps{e}, \; \Theta[x] = \dType \Rightarrow \rho_2[x] = \intoVal{\tyCtx[x]}{\rho_1[x]}{\bot}\right) \\
        \Rightarrow
        s_3 = s_2 \wedge \left(\rType = \dType \Rightarrow v_2 = \intoVal{sty}{v_1}{\bot}\right)^+.
    \end{multline*}

    This lemma formalizes the behavior of an expression during an oversampling step.
    It states that if an expression $e$ is evaluated in two consecutive steps, and all of its
    discrete dependencies (of type \dType) are oversampled in the second step, then the
    computation is stable: the state of the expression does not change ($s_3 = s_2$), and
    its discrete outputs are themselves oversampled versions of their previous values.
\end{lemma}
\begin{skippableproof}
    Let $e$ be an expression in a program \G and typing context \tyCtx, let $sty^+$ be types, let
    $\rType^+$ be reactive types, and let $\Theta$ be a reactive typing context such that:
    \begin{equation*}
        \exprTyped{e}{sty^+} \qquad \exprTypedReact{\Theta}{e}{\rType^+}.
    \end{equation*}
    Let $s_1$, $s_2$, and $s_3$ be states corresponding to $e$, $\rho_1$ and $\rho_2$ be contexts,
    and $v_1^+$ and $v_2^+$ be values such that:
    \begin{equation*}
        \begin{array}{ccc}
            \exprOpsem{\rho_1}{s_1}{e}{s_2}{v_1^+}        &             & \exprOpsem{\rho_2}{s_2}{e}{s_3}{v_2^+}            \\
            \forall x \in \deps{e}, \; \Theta[x] = \dType & \Rightarrow & \rho_2[x] = \intoVal{\tyCtx[x]}{\rho_1[x]}{\bot},
        \end{array}
    \end{equation*}

    We proceed by structural induction on expressions and induction of component calls (recursive
    components are forbidden).

    \paragraph{Constants}
    Suppose $e$ is the constant expression $c$.
    Then we have $s_3 = s_2 = ()$ (\Cref{eq:semantics:grsync:opsem:init:e:const}), and
    the rule \nameref{rule:semantics:grsync:opsem:step:e:const} ensures that $v_1 = v_2 = c$.

    \paragraph{Identifiers}
    Suppose $e$ is the identifier expression $x$.
    Then we have $s_3 = s_2 = ()$ (\Cref{eq:semantics:grsync:opsem:init:e:ident}), and
    the rule \nameref{rule:semantics:grsync:opsem:step:e:ident} ensures that $v_1 =
        \rho_1[x]$ and $v_2 = \rho_2[x]$. Because $x$ is in \deps{x}
    (\Cref{rule:semantics:grsync:causality:dependencies:deps:e:ident}) and $\rType =
        \Theta[x]$ by the rule \nameref{rule:semantics:grsync:react_typing:e:ident}, if
    $\Theta[x] = \dType$ then we have $\rho_2[x] = \intoVal{\tyCtx[x]}{\rho_1[x]}
        {\bot}$ and therefore $v_2 = \intoVal{\tyCtx[x]}{v_1}{\bot}$.

    \paragraph{Last memories}
    Suppose $e$ is the memory expression \last{x}.
    Then we have $s_3 = s_2 = ()$ (\Cref{eq:semantics:grsync:opsem:init:e:last}).
    Because $\rType = \cType$ by the rule \nameref{rule:semantics:grsync:react_typing:e:last}, the
    proof is complete.

    \paragraph{Event emissions}
    Suppose $e$ is the event emission \emit{e'}.
    The states $s_1$, $s_2$, and $s_3$ correspond to $e'$ (\Cref{eq:semantics:grsync:opsem:init:e:emit}).
    %
    The rule \nameref{rule:semantics:grsync:opsem:step:e:emit:value} ensures there exist values
    $v_1'$ and $v_2'$ such that:
    \begin{align*}
        \exprOpsem{\rho_1}{s_1}{e'}{s_2}{v'_1} &  & v_1 & = \some{v'_1}  \\
        \exprOpsem{\rho_2}{s_2}{e'}{s_3}{v'_2} &  & v_2 & = \some{v'_2}.
    \end{align*}
    %
    The rule \nameref{rule:grust:typing:grsync:expr:emit} ensures that there exist a signal type
    $ty$ such that $sty^+ = \mayB{ty}$ and \exprTyped{e'}{ty}.
    %
    For all $x$ in \deps{e'}, we have $x$ in \deps{\emit{e'}}
    (\Cref{rule:semantics:grsync:causality:dependencies:deps:e:emit}), thus if $\Theta[x] = \dType$
    then we have $\rho_2[x] = \intoVal{\tyCtx[x]}{\rho_1[x]}{\bot}$.
    %
    The structural induction hypothesis ensures that $s_3 = s_2$.

    Because $\rType = \cType$ by the rule \nameref{rule:semantics:grsync:react_typing:e:emit}, the
    proof is complete.

    \paragraph{Operations}
    Suppose $e$ is the operation expression $op(e^{\prime+})$.
    For each expression $e'$, there exist states $s_1'$, $s_2'$, and $s_3'$
    corresponding to $e'$, such that $s_1 = (s_1^{\prime+})$, $s_2 = (s_2^{\prime+})$, and
    $s_3 = (s_3^{\prime+})$ (\Cref{eq:semantics:grsync:opsem:init:e:op}).
    %
    The rule \nameref{rule:semantics:grsync:opsem:step:e:op:value} ensures there exist values
    $v_1'$ and $v_2'$ such that:
    \begin{align*}
        \left(\exprOpsem{\rho_1}{s'_1}{e'}{s'_2}{v'_1}\right)^+ &  & v_1 & = op(v_1^{\prime+})  \\
        \left(\exprOpsem{\rho_2}{s'_2}{e'}{s'_3}{v'_2}\right)^+ &  & v_2 & = op(v_2^{\prime+}).
    \end{align*}
    %
    The rule \nameref{rule:grust:typing:grsync:expr:op} ensures that there exist signal types
    $ty^+$ such that $sty^+ = ty^+$, and for each $e'$ there exist a signal type $ty'$
    such that we have \exprTyped{e'}{ty'}.
    %
    For all $x$ in \deps{e'}, we have $x$ in \deps{op(e^{\prime+})} (\Cref{rule:semantics%
        :grsync:causality:dependencies:deps:e:op}), thus if $\Theta[x] = \dType$ then we
    have $\rho_2[x] = \intoVal{\tyCtx[x]}{\rho_1[x]}{\bot}$.

    \begin{itemize}
        \item Suppose that for all expression $e'$, we have \exprTypedReact{\Theta}{e'}
              {\dType}. The rule \nameref{rule:semantics:grsync:react_typing:e:op:%
                  value} ensures that \exprTypedReact{\Theta}{op(e^{\prime+})}
              {\dType}.
              %
              The structural induction hypothesis ensures that $s'_3 = s'_2$ and $v'_2 =
                  \intoVal{ty'}{v'_1}{\bot}= v'_1$.
              %
              Therefore:
              \begin{equation*}
                  \left(v_2 = op(v_2^{\prime+}) = op(v_1^{\prime+}) = v_1 = \intoVal{ty}{v_1}{\bot}\right)^+.
              \end{equation*}

        \item Suppose that there exists an expression $e'$ such that \exprTypedReact
              {\Theta}{e'}{\cType}. Rule \nameref{rule:semantics:grsync:%
                  react_typing:e:op:time} ensures that $\rType = \cType$.
              %
              The structural induction hypothesis ensures that $s'_3 = s'_2$, which
              ensures that $s_3 = s_2$.
    \end{itemize}
    %
    Therefore we have $s_3 = s_2$, and if $\rType = \dType$ then $(v_2 = \intoVal{ty}{v_1}
        {\bot})^+$.

    \paragraph{Component applications}
    Suppose $e$ is the component application
    $f(e^{\prime+})$. For each expression $e'$, there exist states $s_1'$, $s_2'$, and
    $s_3'$ corresponding to $e'$, and there exist states $s^f_1$, $s^f_2$, and $s^f_3$
    corresponding to $f$ such that $s_1 = (s^f_1, s_1^{\prime+})$, $s_2 = (s^f_2,
        s_2^{\prime+})$, and $s_3 = (s^f_3, s_3^{\prime+})$
    (\Cref{eq:semantics:grsync:opsem:init:e:app}).
    %
    The rule \nameref{rule:semantics:grsync:opsem:step:e:app:value} ensures there exist values
    $v_1'$ and $v_2'$ such that:
    \\\noindent\scalebox{0.8}{\begin{minipage}{1.2\textwidth}
            \begin{equation*}
                \begin{array}{llcll}
                     & \left(\exprOpsem{\rho_1}{s'_1}{e'}{s'_2}{v'_1}\right)^+ & \quad &
                     & \compOpsem{s^f_1}{v_1^{\prime+}}{f}{s^f_2}{v_1^+}                 \\
                     & \left(\exprOpsem{\rho_2}{s'_2}{e'}{s'_3}{v'_2}\right)^+ & \quad &
                     & \compOpsem{s^f_2}{v_2^{\prime+}}{f}{s^f_3}{v_2^+}.
                \end{array}
            \end{equation*}
        \end{minipage}}\\\noindent
    %
    The rule \nameref{rule:grust:typing:grsync:expr:compapp} ensures that for each $e'$ there
    exist a type $sty'$ such that we have \exprTyped{e'}{sty'}.
    %
    For all $x$ in \deps{e'}, we have $x$ in\deps{f(e^{\prime+})} (\Cref{rule:semantics:%
        grsync:causality:dependencies:deps:e:app}), thus if $\Theta[x] = \dType$ then
    we have $\rho_2[x] = \intoVal{\tyCtx[x]}{\rho_1[x]}{\bot}$.

    The rule \nameref{rule:semantics:grsync:react_typing:e:app} ensures that for each $e'$
    there exist a reactive type $\rType'$ such that \exprTypedReact{\Theta}{e}
    {\rType'} and \compTypedReact{f}{\rType^{\prime+}}{\rType^+}.
    %
    The structural induction hypothesis ensures that $s'_3 = s'_2$, and if $\rType' =
        \dType$ then $v'_2 = \intoVal{sty'}{v'_1}{\bot}$.

    By induction on component calls, applying \Cref{thm:semantics:grsync:theorems:%
        oversampling}, we have $s_3 = s_2$, and if $\rType = \dType$ then $(v_2 = \intoVal
        {sty}{v_1}{\bot})^+$.

    \setcounter{paragraph}{0}
\end{skippableproof}

\begin{lemma}[Oversampling of \GRSync reactive patterns]
    \label{lem:semantics:grsync:theorems:oversampling:pat}
    Let \guarded{pat^+}{e_g} be a guarded pattern in a program \G and typing context \tyCtx, let
    $sty^+$ be types, let $\rType^+$ be reactive types, and let $\Theta$ be a reactive typing
    context such that:
    \begin{equation*}
        \patTyped{pat^+}{sty^+} \qquad \exprTyped{e_g}{\kwf{bool}} \qquad \patTypedReact{\Theta}{\guarded{pat^+}{e_g}}{\rType^+}.
    \end{equation*}
    For all contexts $\rho_1$, $\rho'_1$, $\rho_2$ and $\rho'_2$, all values $v_1^+$ and $v_2^+$,
    and all booleans $b_1$ and $b_2$ such that:
    \begin{equation*}
        \patOpsem{\rho_1}{\guarded{pat^+ = v_1^+}{e_g}}{b_1}{\rho'_1} \qquad
        \patOpsem{\rho_2}{\guarded{pat^+ = v_2^+}{e_g}}{b_2}{\rho'_2},
    \end{equation*}
    we have:
    \begin{multline*}
        \left(\rType = \dType \Rightarrow v_2 = \intoVal{sty}{v_1}{\bot}\right)^+ \Rightarrow \\
        \left(\forall x \in \deps{e_g} \backslash \defs{pat^+}, \; \Theta[x] = \dType \Rightarrow \rho_2[x] = \intoVal{\tyCtx[x]}{\rho_1[x]}{\bot}\right) \Rightarrow \\
        b_2 = b_1 \wedge \left(\forall y \in \defs{pat^+}, \; \Theta[y] = \dType \Rightarrow \rho'_2[y] = \intoVal{\tyCtx[y]}{\rho'_1[y]}{\bot}\right).
    \end{multline*}

    This lemma extends the oversampling property to guarded patterns. It states that if a
    guarded pattern is evaluated in two consecutive steps where its discrete inputs are
    oversampled---both the values being matched and the dependencies of the guard
    expression---then the outcome is stable: the boolean result of the match remains the
    same ($b_2 = b_1$), and any variables bound by the pattern are themselves correctly
    oversampled.
\end{lemma}
\begin{skippableproof}
    Let \guarded{pat^+}{e_g} be a guarded pattern in a program \G and typing context \tyCtx, let
    $sty^+$ be types, let $\rType^+$ be reactive types, and let $\Theta$ be a reactive typing
    context such that:
    \begin{equation*}
        \patTyped{pat^+}{sty^+} \qquad \exprTyped{e_g}{\kwf{bool}} \qquad \patTypedReact{\Theta}{\guarded{pat^+}{e_g}}{\rType^+}.
    \end{equation*}
    Let $\rho_1$, $\rho'_1$, $\rho_2$ and $\rho'_2$ be contexts, let $v_1^+$ and $v_2^+$ be values,
    and let $b_1$ and $b_2$ be booleans such that:
    \begin{equation*}
        \begin{array}{c}
            \patOpsem{\rho_1}{\guarded{pat^+ = v_1^+}{e_g}}{b_1}{\rho'_1} \qquad \patOpsem{\rho_2}{\guarded{pat^+ = v_2^+}{e_g}}{b_2}{\rho'_2} \\
            \left(\rType = \dType \Rightarrow v_2 = \intoVal{sty}{v_1}{\bot}\right)^+                                                          \\
            \forall x \in \deps{e_g} \backslash \defs{pat^+}, \; \Theta[x] = \dType \Rightarrow \rho_2[x] = \intoVal{\tyCtx[x]}{\rho_1[x]}{\bot}.
        \end{array}
    \end{equation*}

    We proceed by structural induction on patterns.

    \paragraph{Guarded patterns}
    The rule \nameref{rule:semantics:grsync:opsem:step:pats:bool}
    ensures there exist booleans $b_{g1}$ and $b_{g2}$, for each pattern $pat$ booleans
    $b'_1$ and $b'_2$, and for each pattern $pat$ contexts $\rho^{\prime\prime}_1$ and
    $\rho^{\prime\prime}_2$ such that:
    \\\noindent\scalebox{0.8}{\begin{minipage}{1.2\textwidth}
            \begin{equation*}
                \begin{array}{llcll}
                     & \left(\patOpsem{\rho_1}{pat = v_1}{b'_1}{\rho^{\prime\prime}_1}\right)^+                & \quad &
                     & \left(\patOpsem{\rho_2}{pat = v_2}{b'_2}{\rho^{\prime\prime}_2}\right)^+                          \\
                     & \exprOpsem{\rho_1 \ctxconcat \bigCtxconcat \rho^{\prime\prime+}_1}{()}{e_g}{()}{b_{g1}} & \quad &
                     & \exprOpsem{\rho_2 \ctxconcat \bigCtxconcat \rho^{\prime\prime+}_2}{()}{e_g}{()}{b_{g2}}           \\
                     & b_1 = b_{g1} \kwf{\&\&} (\kwf{\&\&} b'_1)^+                                             & \quad &
                     & b_2 = b_{g2} \kwf{\&\&} (\kwf{\&\&} b'_2)^+                                                       \\
                     & \rho'_1 = \bigCtxconcat \rho^{\prime\prime +}_1                                         & \quad &
                     & \rho'_2 = \bigCtxconcat \rho^{\prime\prime +}_2.
                \end{array}
            \end{equation*}
        \end{minipage}}\\\noindent
    %
    The rule \nameref{rule:semantics:grsync:react_typing:pats} ensures that we have
    \exprTypedReact{\Theta}{e_g}{\dType}, and for each $pat$ we have \patTypedReact
    {\Theta}{pat}{\rType}.
    %
    Therefore for each $pat$ we have $b'_2 = b'_1$, and for all $y$ in \defs{pat} such
    that $\Theta[y] = \dType$ we have $\rho^{\prime\prime}_2[y] = \intoVal{\tyCtx[y]}
        {\rho^{\prime\prime}_1[y]}{\bot}$.
    %
    For all $y$ in \defs{pat^+}, there exists a pattern $pat$ such that we have $y$ in
    \defs{pat} (\Cref{rule:semantics:grsync:causality:defs:pat}), thus if $\Theta[y] =
        \dType$ then we have $\bigCtxconcat \rho^{\prime\prime +}_2[y] = \intoVal
        {\tyCtx[x]}{\bigCtxconcat \rho^{\prime\prime +}_1[y]}{\bot}$.
    %
    Therefore, for all $x$ in \deps{e_g}, we have
    $\bigCtxconcat (\rho_2 \ctxconcat \rho^{\prime\prime +}_2)[x] = \intoVal{\tyCtx[x]}
        {(\rho_1 \ctxconcat \bigCtxconcat \rho^{\prime\prime +}_1)[x]}{\bot}$.
    %
    The application of \Cref{lem:semantics:grsync:theorems:oversampling:e} ensures that
    $b_{g2} = \intoVal{\kwf{bool}}{b_{g1}}{\bot} = b_{g1}$.
    %
    It implies that $b_1 = b_{g1} \kwf{\&\&} (\kwf{\&\&} b'_1)^+ = b_{g2} \kwf{\&\&}
        (\kwf{\&\&} b'_2)^+ = b_2$.

    \paragraph{Constants}
    Suppose \guarded{pat^+}{e_g} is the constant pattern $c$.
    Then the rule \nameref{rule:semantics:grsync:opsem:step:pat:const} ensures that $b_1 =
        (c \kwf{==} v_1)$ and $b_2 = (c \kwf{==} v_2)$.
    %
    The rule \nameref{rule:semantics:grsync:react_typing:pat:const} ensures that
    $\rType = \dType$, thus $v_2 = \intoVal{sty}{v_1}{\bot}$.
    %
    The rule \nameref{rule:grust:typing:grsync:pat:const} ensures that $c$ is of signal type,
    thus we have $v_2 = v_1$, which implies $b_2 = b_1$.
    %
    \Cref{rule:semantics:grsync:causality:defs:pat:const} states that \defs{c} is empty.

    \paragraph{Identifiers}
    Suppose \guarded{pat^+}{e_g} is the identifier pattern $x$.
    Then the rule \nameref{rule:semantics:grsync:opsem:step:pat:ident} ensures that $b_1 =
        = b_2 = \true$, $\rho'_1[x] = v_1$, and $\rho'_2[x] = v_2$.
    %
    The rule \nameref{rule:grust:typing:grsync:pat:id} ensures that $sty = \tyCtx[x]$.
    %
    \Cref{rule:semantics:grsync:causality:defs:pat:ident} states that $\defs{x} = \{x\}$.
    %
    The rule \nameref{rule:semantics:grsync:react_typing:pat:ident} ensures that
    $\rType = \Theta[x]$, thus if $\Theta[x] = \dType$ then:
    \begin{equation*}
        \rho'_2[x] = v_2 = \intoVal{sty}{v_1}{\bot} = \intoVal{\tyCtx[x]}{\rho'_1[x]}{\bot}.
    \end{equation*}

    \paragraph{Wildcards}
    Suppose \guarded{pat^+}{e_g} is the wildcard pattern $\_$.
    Then the rule \nameref{rule:semantics:grsync:opsem:step:pat:wild} ensures that $b_1 =
        = b_2 = \true$.
    %
    \Cref{rule:semantics:grsync:causality:defs:pat:wild} states that \defs{x} is empty.

    \setcounter{paragraph}{0}
\end{skippableproof}

\begin{lemma}[Oversampling of \GRSync reactive event patterns]
    \label{lem:semantics:grsync:theorems:oversampling:epat}
    Let \guarded{epat^+}{e_g} be a guarded event pattern in a program \G and typing
    context \tyCtx, and let $\Theta$ be a reactive typing context such that:
    \begin{equation*}
        \epatTyped{epat^+} \qquad \exprTyped{e_g}{\kwf{bool}} \qquad \epatTypedReact{\Theta}{\guarded{epat^+}{e_g}}.
    \end{equation*}
    For all contexts $\rho_1$, $\rho'_1$, $\rho_2$ and $\rho'_2$, all states $s_1$, $s_2$ and $s_3$
    corresponding to \guarded{epat^+}{e_g}, and all booleans $b_1$ and $b_2$ such that:
    \begin{equation*}
        \epatOpsem{\rho_1}{s_1}{\guarded{epat^+}{e_g}}{s_2}{b_1}{\rho'_1} \qquad
        \epatOpsem{\rho_2}{s_2}{\guarded{epat^+}{e_g}}{s_3}{b_2}{\rho'_2},
    \end{equation*}
    we have:
    \begin{multline*}
        \left(\forall x \in \deps{epat^+}, \; \Theta[x] = \dType \Rightarrow \rho_2[x] = \intoVal{\tyCtx[x]}{\rho_1[x]}{\bot}\right) \Rightarrow
        s_3 = s_2 \wedge b_2 = \false.
    \end{multline*}

    This lemma addresses the behavior of event patterns during an oversampling step. It asserts that
    if all discrete dependencies of an event pattern are oversampled, then the pattern will not
    detect a new event ($b_2 = \false$) and its internal state will remain unchanged ($s_3 = s_2$).
    This is a critical property. The reactive type system
    (\Cref{sec:semantics:grsync:react_typing}) ensures that event detection cannot depend on the
    \timeop operator, guaranteeing that events are triggered only by genuine value changes. This
    lemma confirms that, as a result, an oversampling step does not cause spurious event
    activations.
\end{lemma}
\begin{skippableproof}
    Let \guarded{epat^+}{e_g} be a guarded event pattern in a program \G and typing
    context \tyCtx, and let $\Theta$ be a reactive typing context such that:
    \begin{equation*}
        \epatTyped{epat^+} \qquad \exprTyped{e_g}{\kwf{bool}} \qquad \epatTypedReact{\Theta}{\guarded{epat^+}{e_g}}.
    \end{equation*}
    Let $\rho_1$, $\rho'_1$, $\rho_2$ and $\rho'_2$ be contexts, let $s_1$, $s_2$ and $s_3$ be
    states corresponding to \guarded{epat^+}{e_g}, and let $b_1$ and $b_2$ be booleans such that:
    \begin{equation*}
        \begin{array}{c}
            \epatOpsem{\rho_1}{s_1}{\guarded{epat^+}{e_g}}{s_2}{b_1}{\rho'_1} \qquad \epatOpsem{\rho_2}{s_2}{\guarded{epat^+}{e_g}}{s_3}{b_2}{\rho'_2} \\
            \forall x \in \deps{epat^+}, \; \Theta[x] = \dType \Rightarrow \rho_2[x] = \intoVal{\tyCtx[x]}{\rho_1[x]}{\bot}                            \\
            \forall x \in \deps{e_g} \backslash \defs{pat^+}, \; \Theta[x] = \dType \Rightarrow \rho_2[x] = \intoVal{\tyCtx[x]}{\rho_1[x]}{\bot}.
        \end{array}
    \end{equation*}

    We proceed by structural induction on event patterns.

    \paragraph{Guarded event patterns}
    The rule \nameref{rule:semantics:grsync:opsem:step:epats:bool} ensures there exist booleans
    $b_{g1}$ and $b_{g2}$, for each event pattern $epat$ booleans $b'_1$ and $b'_2$, states $s'_1$,
    $s'_2$ and $s'_3$, and contexts $\rho^{\prime\prime}_1$ and $\rho^{\prime\prime}_2$ such that:
    \\\noindent\scalebox{0.8}{\begin{minipage}{1.2\textwidth}
            \begin{equation*}
                \begin{array}{llcll}
                     & \left(\epatOpsem{\rho_1}{s'_1}{epat}{s'_2}{b'_1}{\rho^{\prime\prime}_1}\right)^+        & \quad &
                     & \left(\epatOpsem{\rho_2}{s'_2}{epat}{s'_3}{b'_2}{\rho^{\prime\prime}_2}\right)^+                  \\
                     & \exprOpsem{\rho_1 \ctxconcat \bigCtxconcat \rho^{\prime\prime+}_1}{()}{e_g}{()}{b_{g1}} & \quad &
                     & \exprOpsem{\rho_2 \ctxconcat \bigCtxconcat \rho^{\prime\prime+}_2}{()}{e_g}{()}{b_{g2}}           \\
                     & s_1 = (s'_1)^+ \qquad s_2 = (s'_2)^+                                                    & \quad &
                     & s_3 = (s'_3)^+                                                                                    \\
                     & b_1 = b_{g1} \kwf{\&\&} (\kwf{\&\&} b'_1)^+                                             & \quad &
                     & b_2 = b_{g2} \kwf{\&\&} (\kwf{\&\&} b'_2)^+                                                       \\
                     & \rho'_1 = \bigCtxconcat \rho^{\prime\prime +}_1                                         & \quad &
                     & \rho'_2 = \bigCtxconcat \rho^{\prime\prime +}_2.
                \end{array}
            \end{equation*}
        \end{minipage}}\\\noindent
    %
    The rule \nameref{rule:semantics:grsync:react_typing:epats} ensures that for each
    $epat$ we have \epatTypedReact{\Theta}{epat}.
    %
    \Cref{rule:semantics:grsync:causality:dependencies:deps:epat} ensures that for each
    $epat$ and for all $x$ in \deps{epat} we have $x$ in \deps{epat^+}, and thus if
    $\Theta[x] = \dType$ then $\rho_2[x] = \intoVal{\tyCtx[x]}{\rho_1[x]}{\bot}$.
    %
    Therefore for each $epat$ we have $s'_3 = s'_2$, $b'_2 = \false$, which implies that:
    \begin{align*}
        b_2 & = b_{g2} \kwf{\&\&} (\kwf{\&\&} b'_2)^+ = b_{g2} \kwf{\&\&} \false = \false \\
        s_3 & = (s'_3)^+ = (s'_2)^+ = s_2.
    \end{align*}

    \paragraph{Event detections}
    Suppose \guarded{epat^+}{e_g} is the event detection pattern
    \detect{x}{y}.
    %
    \Cref{rule:semantics:grsync:causality:dependencies:deps:epat:event} states that
    $\deps{\detect{x}{y}} = \{y\}$, and the rule \nameref{rule:semantics:grsync:react_typ%
        ing:epat:event} ensures that $\Theta[y] = \dType$, so that $\rho_2[y] = \intoVal
        {\tyCtx[y]}{\rho_1[y]}{\bot}$.
    %
    The rule \nameref{rule:grust:typing:grsync:epat:event} ensures that $\tyCtx[y]$ is of the
    form \mayB{ty} which implies that $\rho_2[y] = \none$ and
    \nameref{rule:semantics:grsync:opsem:step:epat:event:none} ensure that $b_2 = \false$.

    \paragraph{Rising edges}
    Suppose \guarded{epat^+}{e_g} is the rising edge $e$.
    %
    \Cref{rule:semantics:grsync:causality:dependencies:deps:epat:edge} states that
    its dependency set equals the dependency set of the expression $e$, and the rule
    \nameref{rule:semantics:grsync:react_typing:epat:edge} ensures that \exprTypedReact
    {\Theta}{e}{\dType}.
    %
    The rule \nameref{rule:semantics:grsync:opsem:step:epat:edge:bool} ensures there exist
    a value $v_m$, states $s'_1$, $s'_2$ and $s'_3$ corresponding to the
    expression $e$, and booleans $b'_1$ and $b'_2$ such that:
    \\\noindent\scalebox{0.8}{\begin{minipage}{1.2\textwidth}
            \begin{equation*}
                \begin{array}{llcll}
                     & \exprOpsem{\rho_1}{s'_1}{e}{s'_2}{b'_1}     & \quad &
                     & \exprOpsem{\rho_2}{s'_2}{e}{s'_3}{b'_2}               \\
                     & s_1 = (v_m, s'_1) \qquad s_2 = (b'_1, s'_2) & \quad &
                     & s_3 = (b'_2, s'_3)                                    \\
                     & b_1 = \kwf{!}v_m \kwf{\&\&} b'_1            & \quad &
                     & b_2 = \kwf{!}b'_1 \kwf{\&\&} b'_2                     \\
                     & \rho'_1 = \emptyrho                         & \quad &
                     & \rho'_2 = \emptyrho.
                \end{array}
            \end{equation*}
        \end{minipage}}\\\noindent
    %
    The rule \nameref{rule:grust:typing:grsync:epat:edge} ensures that \exprTyped{e}{\kwf{bool}}.
    %
    The oversampling \Cref{lem:semantics:grsync:theorems:oversampling:e} on expressions
    ensures that we have  $s'_3 = s'_2$ and $b'_2 = \intoVal{\kwf{bool}}{b'_1}{\bot} =
        b'_1$, which implies that:
    \begin{align*}
        b_2 & = \kwf{!}b'_1 \kwf{\&\&} b'_2 = \kwf{!}b'_1 \kwf{\&\&} b'_1 = \false \\
        s_3 & = (b'_2, s'_3) = (b'_1, s'_2) = s_2.
    \end{align*}

    \setcounter{paragraph}{0}
\end{skippableproof}

\begin{lemma}[Oversampling of \GRSync reactive equations]
    \label{lem:semantics:grsync:theorems:oversampling:eq}
    Let $eq^+$ be a causal set of equations in a program \G and typing context \tyCtx, and let
    $\Theta$ be a reactive typing context such that:
    \begin{equation*}
        \eqTyped{eq^+} \qquad \eqTypedReact{\Theta}{eq^+}.
    \end{equation*}
    For all contexts $\rho_1$, $\rho'_1$, $\rho_2$ and $\rho'_2$, and all states $s_1$, $s_2$ and
    $s_3$ corresponding to $eq^+$ such that:
    \begin{equation*}
        \eqOpsem{\rho_1}{s_1}{eq^+}{s_2}{\rho'_1} \qquad
        \eqOpsem{\rho_2}{s_2}{eq^+}{s_3}{\rho'_2},
    \end{equation*}
    we have:
    \begin{multline*}
        \left(\forall x \in \deps{eq^+}, \; \Theta[x] = \dType \Rightarrow \rho_2[x] = \intoVal{\tyCtx[x]}{\rho_1[x]}{\bot}\right) \Rightarrow \\
        s_3 = s_2 \wedge \left(\forall y \in \defs{eq^+}, \; \Theta[y] = \dType \Rightarrow \rho'_2[y] = \intoVal{\tyCtx[y]}{\rho'_1[y]}{\bot}\right).
    \end{multline*}

    This lemma lifts the oversampling property from individual expressions and patterns to a
    causal set of reactive equations. It establishes that if a block of equations is
    evaluated in two consecutive steps, and all of its discrete dependencies are oversampled
    in the second step, then the entire computation remains stable. Specifically, the state
    managed by the equations does not change ($s_3 = s_2$), and any values of a discrete
    type defined within the equations are themselves correctly oversampled.
\end{lemma}
\begin{skippableproof}
    Let $eq^+$ be a causal set of equations in a program \G and typing context \tyCtx, and let
    $\Theta$ be a reactive typing context such that:
    \begin{equation*}
        \eqTyped{eq^+} \qquad \eqTypedReact{\Theta}{eq^+}.
    \end{equation*}
    Let $\rho_1$, $\rho'_1$, $\rho_2$ and $\rho'_2$ be contexts, $s_1$, $s_2$ and $s_3$ be states
    corresponding to $eq^+$ such that:
    \begin{align*}
        \eqOpsem{\rho_1}{s_1}{eq^+}{s_2}{\rho'_1} \qquad \eqOpsem{\rho_2}{s_2}{eq^+}{s_3}{\rho'_2} \\
        \forall x \in \deps{eq^+}, \; \Theta[x] = \dType \Rightarrow \rho_2[x] = \intoVal{\tyCtx[x]}{\rho_1[x]}{\bot}.
    \end{align*}

    We proceed by structural induction on \GRSync equations.

    \paragraph{Overview of the proof strategy}
    The difficult points of the proof are the following.
    \begin{itemize}
        \item For blocks of mutually recursive equations, we use rank induction
              (\Cref{def:semantics:grsync:causality:induction:rank}) to prove iteratively the
              oversampling of identifiers following the rank
              (\Cref{def:semantics:grsync:causality:rank}).
        \item For \kwf{match} equations, the main difficulty is to prove that the control flow
              remains stable, \ie that the same branch is chosen in both steps. This is achieved by
              applying the oversampling lemma for patterns
              (\Cref{lem:semantics:grsync:theorems:oversampling:pat}). This lemma ensures the
              boolean outcome of each pattern match is unchanged, thus guaranteeing the same branch
              is executed.
        \item For \kwf{when} equations, the challenge is to show that an oversampling step does not
              trigger the stateful computations within a branch. The proof uses the oversampling
              lemma for event patterns (\Cref{lem:semantics:grsync:theorems:oversampling:epat}) to
              demonstrate that no event is detected. This forces the execution to follow the
              non-activating path, which by definition preserves the state and produces the correct
              oversampled outputs.
    \end{itemize}

    \paragraph{Block of equations}
    \Cref{rule:semantics:grsync:causality:dependencies:deps:eq,%
        rule:semantics:grsync:causality:defs:eq} ensure that:
    \begin{align*}
        \deps{eq^+} & = \left(\bigcup \deps{eq}^+\right) \backslash \defs{eq^+} \\
        \defs{eq^+} & = \bigsqcup \defs{eq}^+.
    \end{align*}
    %
    The rule \nameref{rule:semantics:grsync:react_typing:eqs} ensures that for each
    individual equation $eq$ we have \eqTypedReact{\Theta}{eq}.
    %
    The rule \nameref{rule:semantics:grsync:opsem:step:eqs} ensures that for each
    individual equation $eq$ there exist states $s^{eq}_1$, $s^{eq}_2$ and $s^{eq}_3$
    corresponding to $eq$ such that $s_1 = ((s_1^{eq})^+)$, $s_2 = ((s_2^{eq})^+)$ and
    $s_3 = ((s_3^{eq})^+)$.

    The causality of $eq^+$ allows to use the rank induction
    (\Cref{def:semantics:grsync:causality:induction:rank}) to determine for all
    $x \in \defs{eq^+}$ the property \P{x}: if the equation $eq_x$ defines $x$, then:
    \begin{align*}
        s_3^{eq_x} = s_2^{eq_x} \\
        \Theta[x] = \dType \Rightarrow \rho'_2[x] = \intoVal{\tyCtx[x]}{\rho'_1[x]}{\bot}.
    \end{align*}

    \begin{itemize}
        \item \textit{Base case:} Let $x \in \defs{eq^+}$ such that $\rank{eq^+}{x} = 0$.

              Therefore, there exists an equation $eq_x$ that defines $x$ and that has no
              dependency with \defs{eq^+}. By induction hypothesis on equations structure,
              we have:
              \begin{align*}
                   & s_3^{eq_x} = s_2^{eq_x}                                                                                          \\
                   & \forall z \in \defs{eq_x}, \; \Theta[z] = \dType \Rightarrow \rho'_2[z] = \intoVal{\tyCtx[z]}{\rho'_1[z]}{\bot},
              \end{align*}
              in particular:
              \begin{equation*}
                  \Theta[x] = \dType \Rightarrow \rho'_2[x] = \intoVal{\tyCtx[x]}{\rho'_1[x]}{\bot},
              \end{equation*}

        \item \textit{Inductive step:} Let $k \in \N$.
              We assume that for all $y \in \defs{eq^+}$ such that $\rank{eq^+}{y} \le k$
              we have:
              \begin{align*}
                   & \exists eq_y \in \{eq^+\}, \; y \in \defs{eq_y} \wedge s_3^{eq_y} = s_2^{eq_y}     \\
                   & \Theta[y] = \dType \Rightarrow \rho'_2[y] = \intoVal{\tyCtx[y]}{\rho'_1[y]}{\bot}.
              \end{align*}

              Let $x \in \defs{eq^+}$ such that $\rank{eq^+}{x} = k+1$. For all
              $y \in \defs{eq^+}$ such that $x$ depends on $y$ (\ie \depOrd{eq^+}{x}{y}),
              \Cref{lem:semantics:grsync:causality:rank:decrease} states that
              $\rank{eq^+}{x} > \rank{eq^+}{y}$, which implies that
              $\rank{eq^+}{y} \le k$. By rank induction hypothesis, for all
              $y \in \defs{eq^+}$ such that \depOrd{eq^+}{x}{y}, we have:
              \begin{equation*}
                  \Theta[y] = \dType \Rightarrow \rho'_2[y] = \intoVal{\tyCtx[y]}{\rho'_1[y]}{\bot}.
              \end{equation*}

              Let $eq_x$ be the equation that defines $x$. Then for all
              $y \in \deps{eq_x}$, we have:
              \begin{equation*}
                  \Theta[y] = \dType \Rightarrow \rho'_2[y] = \intoVal{\tyCtx[y]}{\rho'_1[y]}{\bot}.
              \end{equation*}
              By induction hypothesis on equations, for all $z$ in \defs{eq_z} we have:
              \begin{align*}
                   & s_3^{eq_x} = s_2^{eq_x}                                                                                          \\
                   & \forall z \in \defs{eq_x}, \; \Theta[z] = \dType \Rightarrow \rho'_2[z] = \intoVal{\tyCtx[z]}{\rho'_1[z]}{\bot},
              \end{align*}
              in particular:
              \begin{equation*}
                  \Theta[x] = \dType \Rightarrow \rho'_2[x] = \intoVal{\tyCtx[x]}{\rho'_1[x]}{\bot},
              \end{equation*}
    \end{itemize}
    Therefore, the property \P{x} holds for all $x \in \defs{eq^+}$.

    It implies that:
    \begin{align*}
         & s_3 = (s_3^{eq+}) = (s_2^{eq+}) = s_2                                                                            \\
         & \forall x \in \defs{eq^+}, \; \Theta[x] = \dType \Rightarrow \rho'_2[x] = \intoVal{\tyCtx[x]}{\rho'_1[x]}{\bot}.
    \end{align*}

    \paragraph{Declarations}
    Suppose that $eq^+$ is the declaration $x^+ = e$.
    %
    \Cref{rule:semantics:grsync:causality:dependencies:deps:eq:decl} ensures that
    $\deps{x^+ = e} = \deps{e}$.
    %
    \Cref{rule:semantics:grsync:causality:defs:eq:decl} ensures that
    $\defs{x^+ = e} = \{x^+\}$.
    %
    The rule \nameref{rule:semantics:grsync:opsem:step:eq:decl} ensures that there exists
    values $v_1^+$ and $v_2^+$ such that:
    \\\noindent\scalebox{0.8}{\begin{minipage}{1.2\textwidth}
            \begin{equation*}
                \begin{array}{llcll}
                     & \exprOpsem{\rho_1}{s_1}{e}{s_2}{v_1^+} & \quad &
                     & \exprOpsem{\rho_2}{s_2}{e}{s_3}{v_2^+}           \\
                     & \left(\rho'_1[x] = v_1\right)^+        & \quad &
                     & \left(\rho'_2[x] = v_2\right)^+ ,
                \end{array}
            \end{equation*}
        \end{minipage}}\\\noindent
    so that for all $x$ in \deps{e} such that $\Theta[x] = \dType$ we have $\rho_2[x] =
        \intoVal{\tyCtx[x]}{\rho_1[x]}{\bot}$.
    %
    The rule \nameref{rule:semantics:grsync:react_typing:eq:decl} ensures that
    \exprTypedReact{\Theta}{e}{(\Theta[x])^+}.
    %
    The oversampling \Cref{lem:semantics:grsync:theorems:oversampling:e} on expressions
    ensures that:
    \begin{align*}
        s_3 = s_2                     &                          \\
        \forall y \in \defs{y^+ = e}, &                          \\
        \Theta[y] = \dType            & \Rightarrow \rho'_2[y] =
        v_2 = \intoVal{\tyCtx[y]}{v_1}{\bot} = \intoVal{\tyCtx[y]}{\rho'_1[y]}{\bot}.
    \end{align*}

    \paragraph{Match equations}
    Suppose that $eq^+$ is the \kwf{match} equation:
    \begin{equation*}
        \match{e}{x^+}{(\branch{pat_k^+}{e_k}{eq_k^+})^+}.
    \end{equation*}
    %
    \Cref{%
        rule:semantics:grsync:causality:dependencies:deps:eq:match,%
        rule:semantics:grsync:causality:defs:eq:match%
    } ensures that:
    \begin{align*}
        \deps{\mathtiny{\matchTab{e}{x^+}{(\branch{pat_k^+}{e_k}{eq_k^+})^+}}} & =
        \deps{e} \cup \bigcup \left(\deps{e_k} \cup \deps{eq_k^+} \backslash \defs{pat_k^+}\right)^+ \\
        \defs{\mathtiny{\matchTab{e}{x^+}{(\branch{pat_k^+}{e_k}{eq_k^+})^+}}} & = \{x^+\}.
    \end{align*}
    %
    The rule \nameref{rule:grust:typing:grsync:eq:match} ensures there exist types $sty^+$ such
    that \exprTyped{e}{sty^+}, for all patterns \guarded{pat_k^+}{e_k} we have \patTyped
    {pat_k^+}{sty^+} and \exprTyped{e_k}{\kwf{bool}}, and for all block of equations
    $eq_k^+$ we have \eqTyped{eq_k^+}.
    %
    The rule \nameref{rule:semantics:grsync:react_typing:eq:match} ensures there exist
    reactive types $\rType^+$ such that:
    \begin{align*}
         & \exprTypedReact{\Theta}{e}{\rType^+}                                   \\
         & \forall k, \; \patTypedReact{\Theta}{\guarded{pat_k^+}{e_k}}{\rType^+} \\
         & \forall k, \; \eqTypedReact{\Theta}{eq_k^+}.
    \end{align*}
    %
    The rule \nameref{rule:semantics:grsync:opsem:step:eq:match:value} ensures that there exists
    states $s_{e1}$, $s_{e2}$ and $s_{e3}$ corresponding to $e$, values $v_1^+$ and
    $v_2^+$, for each pattern \guarded{pat_k^+}{e_k} there exist booleans $b_{k1}$ and
    $b_{k2}$ and contexts $rho_{k1}$ and $rho_{k2}$, identifiers $i_1$ and $i_2$, and for
    each $eq_k^+$ there exist states $s_{k1}$, $s_{k2}$ and $s_{k3}$ corresponding to each
    $eq_k^+$ such that:
    \\\noindent\scalebox{0.8}{\begin{minipage}{1.2\textwidth}
            \begin{equation*}
                \begin{array}{llcll}
                                        & s_1 = (s_{e1}, s_{k1}^+) \qquad s_2 = (s_{e2}, s_{k2}^+)                               & \quad &
                                        & s_3 = (s_{e3}, s_{k3}^+)                                                                         \\
                                        & \exprOpsem{\rho_1}{s_{e1}}{e}{s_{e2}}{v_1^+}                                           & \quad &
                                        & \exprOpsem{\rho_2}{s_{e2}}{e}{s_{e3}}{v_2^+}                                                     \\
                    \forall k,          & \patOpsem{\rho_1}{\guarded{pat_k^+ = v_1^+}{e_k}}{b_{k1}}{\rho_{k1}}                   & \quad &
                    \forall k,          & \patOpsem{\rho_2}{\guarded{pat_k^+ = v_2^+}{e_k}}{b_{k2}}{\rho_{k2}}                             \\
                    \forall k < i_1,    & b_k = \false \wedge b_{i_1} = \true                                                    & \quad &
                    \forall k < i_2,    & b_k = \false \wedge b_{i_2} = \true                                                              \\
                                        & \eqOpsem{\rho_1 \ctxconcat \rho_{{i_1}1}}{s_{{i_1}1}}{eq_{i_1}^+}{s_{{i_1}2}}{\rho'_1} & \quad &
                                        & \eqOpsem{\rho_2 \ctxconcat \rho_{{i_2}2}}{s_{{i_2}2}}{eq_{i_2}^+}{s_{{i_2}3}}{\rho'_2}           \\
                    \forall k \neq i_1, & s_{k2} = s_{k1}                                                                        & \quad &
                    \forall k \neq i_2, & s_{k3} = s_{k2},
                \end{array}
            \end{equation*}
        \end{minipage}}\\\noindent
    so that for all $x$ in \deps{e} such that $\Theta[x] = \dType$ we have $\rho_2[x] =
        \intoVal{\tyCtx[x]}{\rho_1[x]}{\bot}$.
    %
    The oversampling \Cref{lem:semantics:grsync:theorems:oversampling:e} on expressions
    ensures that:
    \begin{align*}
         & s_{e3} = s_{e2}                                                            \\
         & \left(\rType = \dType \Rightarrow v_2 = \intoVal{sty}{v_1}{\bot}\right)^+.
    \end{align*}
    %
    The oversampling \Cref{lem:semantics:grsync:theorems:oversampling:pat} on patterns
    ensures that:
    \begin{align*}
        \forall k, \; & b_{k2} = b_{k1}                                                                                                                      \\
        \forall k, \; & \left(\forall y \in \defs{pat_k^+}, \; \Theta[y] = \dType \Rightarrow \rho_{k2}[y] = \intoVal{\tyCtx[y]}{\rho_{k1}[y]}{\bot}\right),
    \end{align*}
    so that we have $i_1 = i_2 = i$, and for all $x$ in \deps{eq_i^+} such that $Theta[x]
        = \dType$ we have $(\rho_2 \ctxconcat \rho_{i2})[x] = \intoVal{\tyCtx[x]}{(\rho_1
            \ctxconcat \rho_{i1})[x]}{\bot}$.
    %
    The application of the structural induction hypothesis on equations ensures that:
    \begin{align*}
         & s_{i3} = s_{i2}                                                                                                    \\
         & \forall y \in \defs{eq_i^+}, \; \Theta[y] = \dType \Rightarrow \rho'_2[y] = \intoVal{\tyCtx[y]}{\rho'_1[y]}{\bot}.
    \end{align*}
    %
    Moreover, for all $k \neq i$, we have $s_{k3} = s_{k2}$.
    % 
    \Cref{rule:semantics:grsync:causality:defs:eq:match} ensures that $\defs{\kwf{match}}
        = \defs{eq_i^+}$, so that:
    \begin{align*}
         & s_3 = (s_{e3}, s_{k3}^+) = (s_{e2}, s_{k2}^+) = s_2                                                                     \\
         & \forall y \in \defs{\kwf{match}}, \; \Theta[y] = \dType \Rightarrow \rho'_2[y] = \intoVal{\tyCtx[y]}{\rho'_1[y]}{\bot}.
    \end{align*}

    \paragraph{When equations}
    Suppose that $eq^+$ is the \kwf{when} equation:
    \begin{equation*}
        \when{x^+}{x_m^+ = c^+}{(\branch{epat_k^+}{e_k}{eq_k^+})^+}.
    \end{equation*}
    %
    \Cref{%
        rule:semantics:grsync:causality:dependencies:deps:eq:when,%
        rule:semantics:grsync:causality:defs:eq:when%
    } ensures that:
    \begin{align*}
        \deps{\kwf{when}} & =
        \begin{array}{l}
            \bigcup \left(\deps{epat_k^+} \cup \deps{e_k} \cup \deps{eq_k^+} \backslash \defs{epat_k^+}\right)^+ \\
            \qquad \backslash (\{x^+\} \cup \{\last{x_m}^+\})
        \end{array} \\
        \defs{\kwf{when}} & = \{x^+\}.
    \end{align*}
    %
    The rule \nameref{rule:semantics:grsync:react_typing:eq:when} ensures that:
    \begin{align*}
         & \forall x \in \defs{\kwf{when}}, \; \Theta[x] = \dType          \\
         & \forall k, \; \epatTypedReact{\Theta}{\guarded{epat_k^+}{e_k}}.
    \end{align*}
    %
    The rules \nameref{rule:semantics:grsync:opsem:step:eq:when:branch} and
    \nameref{rule:semantics:grsync:opsem:step:eq:when:default} ensure that there exists values
    $v_m^+$, for each event pattern \guarded{epat_k^+}{e_k} there exist states $s^p_{k1}$,
    $s^p_{k2}$ and $s^p_{k3}$ corresponding to \guarded{epat_k^+}{e_k}, booleans $b_{k1}$
    and $b_{k2}$ and contexts $rho_{k1}$ and $rho_{k2}$, and for each $eq_k^+$ there exist
    states $s_{k1}$, $s_{k2}$ and $s_{k3}$ corresponding to each $eq_k^+$ such that:
    \\\noindent\scalebox{0.8}{\begin{minipage}{1.2\textwidth}
            \begin{equation*}
                \begin{array}{llcll}
                    \forall k, & \epatOpsem{\rho_1}{s^p_{k1}}{\guarded{epat_k^+}{e_k}}{s^p_{k2}}{b_{k1}}{\rho_{k1}} & \quad &
                    \forall k, & \epatOpsem{\rho_2}{s^p_{k2}}{\guarded{epat_k^+}{e_k}}{s^p_{k3}}{b_{k2}}{\rho_{k2}}           \\
                               & s_1 = (v_m^+, s^{p+}_{k1}, s_{k1}^+)                                                         \\
                               & s_2 = ((\rho'_1[x_m])^+, s^{p+}_{k2}, s_{k2}^+)                                    & \quad &
                               & s_3 = ((\rho'_2[x_m])^+, s^{p+}_{k3}, s_{k3}^+).
                \end{array}
            \end{equation*}
        \end{minipage}}\\\noindent
    %
    The oversampling \Cref{lem:semantics:grsync:theorems:oversampling:epat} on event
    patterns ensures that for all $k$ we have $s^p_{k3} = s^p_{k2}$ and $b_{k2} = \false$.
    %
    The rule \nameref{rule:semantics:grsync:opsem:step:eq:when:default} ensures that:
    \begin{align*}
         & \left(\rho'_2[x_m] = \rho'_1[x_m]\right)^+                                                                                               \\
         & \left(s_{k3} = s_{k2}\right)^+                                                                                                           \\
         & \rho'_2 = [(\rho'_1[x_m] / \last{x_m})^+] \ctxconcat [\none / x_e \in \{x^+\} \backslash \{x_m^+\}] \ctxconcat [(\rho'_1[x_m] / x_m)^+],
    \end{align*}
    so that $s_3 = s_2$, and for all $y$ in \defs{\kwf{when}} we have:
    \begin{itemize}
        \item if $y$ is in $\{x_m\}$, then its type corresponds to a signal type $ty$
              (\cf the typing rules \nameref{rule:grust:typing:grsync:eq:when}
              and \nameref{rule:grust:typing:grsync:expr:const}) and $\rho'_2[y] = \rho'_1[y]$
              which equals \intoVal{ty}{\rho'_1[y]}{\bot}; and
        \item if $y$ is not in $\{x_m\}$, then $y$ is of event type \mayB{ty} (\cf the
              typing rule \nameref{rule:grust:typing:grsync:eq:when}) and $\rho'_2[y] = \none$
              which equals \intoVal{\mayB{ty}}{\rho'_1[y]}{\bot}.
    \end{itemize}
    %
    Therefore for all $y$ in \defs{\kwf{when}} we have $\rho'_2[y] = \intoVal{\tyCtx[y]}
        {\rho'_1[y]}{\bot}$.

    \setcounter{paragraph}{0}
\end{skippableproof}

\begin{theorem}[Oversampling]
    \label{thm:semantics:grsync:theorems:oversampling}
    Let \component{f}{x_i^+}{x_o^+}{eq^+}{x_m^+ = c^+} be a causal and well-typed component in a
    program \G and typing context \tyCtx, and let $\rType_i^+$ and $\rType_o^+$ be input and
    output reactive types such that:
    \begin{equation*}
        \compTypedReact{f}{\rType_i^+}{\rType_o^+}.
    \end{equation*}
    For all states $s_1$, $s_2$, and $s_3$ corresponding to $f$, and all values $v_{i1}^+$,
    $v_{i2}^+$, $v_{o1}^+$ and $v_{o2}^+$, such that:
    \begin{equation*}
        \compOpsem{s_1}{v_{i1}^+}{f}{s_2}{v_{o1}^+} \qquad \compOpsem{s_2}{v_{i2}^+}{f}{s_3}{v_{o2}^+},
    \end{equation*}
    we have:
    \begin{multline*}
        \left(\forall x_i, \; \rType_i = \dType \Rightarrow v_{i2} = \intoVal{\tyCtx[x_i]}{v_{i1}}{\bot}\right) \\
        \Rightarrow
        s_3 = s_2 \wedge \left(\forall x_o, \; \rType_o = \dType \Rightarrow v_{o2} = \intoVal{\tyCtx[x_o]}{v_{o1}}{\bot}\right).
    \end{multline*}

    This theorem states the oversampling property to the component level. It guarantees that if a
    reactive component is executed on oversampled discrete inputs, its behavior remains stable: its
    internal state does not change ($s_3 = s_2$), and its discrete outputs are correctly
    oversampled. This result is crucial, proving that the behavior of a \GRSync  component is driven
    solely by discrete input signal changes and event occurrences. This independence from the
    execution rate is essential for the safe integration of components into \GRFRP services, which
    are designed to propagate only signal changes and event occurrences.
\end{theorem}
\begin{skippableproof}
    Let \component{f}{x_i^+}{x_o^+}{eq^+}{x_m^+ = c^+} be a causal and well-typed component in a
    program \G and typing context \tyCtx, and let $\rType_i^+$ and $\rType_o^+$ be input and
    output reactive types such that:
    \begin{equation*}
        \compTypedReact{f}{\rType_i^+}{\rType_o^+}.
    \end{equation*}
    Let $s_1$, $s_2$, and $s_3$ be states corresponding to $f$, let $v_{i1}^+$, $v_{i2}^+$,
    $v_{o1}^+$ and $v_{o2}^+$ be values such that:
    \begin{align*}
         & \compOpsem{s_1}{v_{i1}^+}{f}{s_2}{v_{o1}^+} \qquad \compOpsem{s_2}{v_{i2}^+}{f}{s_3}{v_{o2}^+} \\
         & \forall x_i, \; \rType_i = \dType \Rightarrow v_{i2} = \intoVal{\tyCtx[x_i]}{v_{i1}}{\bot}.
    \end{align*}

    The rule \nameref{rule:semantics:grsync:react_typing:comp} ensures that:
    \begin{align*}
        \left(\Theta[x_i] = \rType_i\right)^+       &  &  &
        \left(\Theta[x_o] = \rType_o\right)^+       &  &  &
        \left(\Theta[x_m] = \dType \right)^+        &  &  &
        \left(\Theta[\last{x_m}] = \cType \right)^+ &  &  &
        \eqTypedReact{\Theta}{eq^+}.
    \end{align*}

    The rule \nameref{rule:semantics:grsync:opsem:step:comp} ensures there exist values $v_m^+$,
    states $s'_1$, $s'_2$, and $s'_3$ corresponding to $eq^+$, and contexts $\rho_1$, $\rho_2$,
    $\rho'_1$ and $\rho'_2$ such that:
    \begin{center}
        \scalebox{0.8}{\begin{minipage}{1.2\textwidth}
            \begin{equation*}
                \begin{array}{llcll}
                     & s_1 = (v_m^+, s'_1) \qquad s_2 = (\rho'_1[x_m]^+, s'_2) & \quad &
                     & s_3 = (\rho'_2[x_m]^+, s'_3)                                      \\
                     & \rho_1 = [(v_{i1}/x_i)^+, (v_m/\last{x_m})^+]           & \quad &
                     & \rho_2 = [(v_{i2}/x_i)^+, (\rho'_1[x_m]/\last{x_m})^+]            \\
                     & \eqOpsem{\rho_1}{s'_1}{eq^+}{s'_2}{\rho'_1}             & \quad &
                     & \eqOpsem{\rho_2}{s'_2}{eq^+}{s'_3}{\rho'_2}                       \\
                     & \left(v_{o1} = \rho'_1[x_o]\right)^+                    & \quad &
                     & \left(v_{o2} = \rho'_2[x_o]\right)^+.
                \end{array}
            \end{equation*}
        \end{minipage}}
    \end{center}

    Because $f$ is well-typed, $\deps{eq^+} \subseteq \{x_i^+, \last{x_m}^+\}$ and $\{x_o^+, x_m^+\}
        \subseteq \defs{eq^+}$.
    %
    Therefore, for all $x$ in \deps{eq^+}, if $x$ is in $\{x_i^+\}$ and $\Theta[x] = \dType$
    then we have $\rho_2[x] = v_{i2} = \intoVal{\tyCtx[x_i]}{v_{i1}}{\bot} = \rho_1[x]$,
    otherwise $x$ is in $\{\last{x_m}^+\}$ and $\Theta[x] = \cType \neq \dType$.

    The oversampling \Cref{lem:semantics:grsync:theorems:oversampling:eq} on equations ensures that:
    \begin{align*}
         & s'_3 = s'_2                                                                                                      \\
         & \forall y \in \defs{eq^+}, \; \Theta[y] = \dType \Rightarrow \rho'_2[y] = \intoVal{\tyCtx[y]}{\rho'_1[y]}{\bot}.
    \end{align*}

    The rule \nameref{rule:grust:typing:grsync:eq:init} ensures that, because constants are of signal type,
    only signal typed memories are created.
    %
    Therefore for all initialized memory $x_m$, there exist a type $ty_m$ such that $ty_m =
        \tyCtx[x_m]$. Furthermore, for all $x_m$ we have $\Theta[x_m] = \dType$, thus:
    \begin{equation*}
        \forall x_m, \; \rho'_2[x_m] = \intoVal{ty_m}{\rho'_1[x_m]}{\bot} = \rho'_1[x_m],
    \end{equation*}
    which implies:
    \begin{align*}
         & s_3 = (\rho'_2[x_m]^+, s'_3) = (\rho'_1[x_m]^+, s'_2) = s_2                                                                                            \\
         & \forall x_o, \; \rType_o = \dType \Rightarrow v_{o2} = \rho'_2[x_o] = \intoVal{\tyCtx[x_o]}{\rho'_1[x_o]}{\bot} = \intoVal{\tyCtx[x_o]}{v_{o1}}{\bot}.
    \end{align*}
\end{skippableproof}

\paragraph{Summary}
In conclusion, this section established the oversampling property for reactive \GRSync components.
%
The core result, \Cref{thm:semantics:grsync:theorems:oversampling}, formally guarantees that
executing a reactive component with unchanged discrete inputs---a process called oversampling---%
results in a stable computation.
%
Specifically, the internal state of the component remains unchanged, and its discrete outputs are
correctly oversampled, meaning signals retain their value and events do not occur.
%
The proof was demonstrated by structural induction, with key lemmas confirming that oversampling
does not alter control flow in \kwf{match} statements or trigger stateful computations in
\kwf{when} clauses.
%
This theorem provides a critical justification for the compiler optimization that shifts from a
periodic to an event-driven execution model. It proves that the behavior of a reactive component is
driven solely by meaningful changes in its inputs, ensuring its safe and correct integration into
the asynchronous environment of \GRFRP services.
